\subsection*{先识别固定相对增长率的含义}
指数模型 $N'=rN$ 说的是变化量与当前数量成比例，所以每过同样时间，数量乘上同样倍率。若 $r>0$，翻倍所需时间满足 $\mathrm e^{rT}=2$，得到 $T=\ln2/r$，与当前绝对数量无关。

这解释了早期看似缓慢、之后迅速增加的曲线形状，但形状本身不能证明背后机制。任何固定相对增长率模型，都暗含条件在研究期间保持足够稳定；当资源或可参与群体有限时，这个假设会逐渐失效。

\subsection*{给增长规则加入容量限制}
一个教学修正是 Logistic 方程 $N'=rN(1-N/K)$。数量远小于 $K$ 时，括号接近一，恢复指数近似；数量接近 $K$ 时，相对增长率减小，最终趋向平衡。这里 $K$ 是模型设定的容量，不能不经估计就赋给真实系统。

把增长率看成 $N$ 的函数，它是开口向下的二次式，在 $K/2$ 最大。于是总数量还在增加时，增加速度已经可以开始下降。区分状态与变化率，正是高阶导数章节中的同一个主题。

本节用于理解数学模型，不给实际疫情预测。真实观测还受到报告延迟、检测和行为变化等因素影响。合理的学习顺序是先从假设推导后果，再对照数据检验假设；不能因为一条曲线像指数或 S 形，就认定机制已经确定。
